% =====================================================================
%  BAB II LANDASAN TEORI
% =====================================================================

\chapter{LANDASAN TEORI}

\section{Tampilan Tujuh Segmen}

Tampilan tujuh segmen menyajikan angka menggunakan tujuh batang segmen yang
disusun seperti angka delapan, ditambah satu titik desimal. Setiap digit
dibentuk oleh kombinasi segmen yang menyala, bukan oleh bentuk huruf. Karena
itu pendekatan pengenalan berbasis huruf tidak dapat langsung diterapkan.
Kekhasan tersebut terlihat pada angka ``1'' yang hanya menyalakan dua segmen
tegak dan angka ``8'' yang menyalakan seluruh segmen, sehingga jarak kesamaan
antarkelas tidak seragam. Titik desimal berukuran sangat kecil, umumnya tiga
hingga lima piksel pada citra beresolusi sedang, sehingga rawan hilang pada
proses pengenalan \parencite{gushima2023automatic}.

Tampilan dibedakan menjadi dua jenis berdasarkan sumber cahayanya. Tampilan
\textit{light emitting diode} memancarkan cahaya sendiri, umumnya berwarna
merah, hijau, atau biru, dengan latar gelap. Tampilan \textit{liquid crystal
display} pasif tidak memancarkan cahaya dan bergantung pada cahaya sekitar,
sehingga kontrasnya rendah dan sudut pandangnya terbatas. Kedua jenis tersebut
memerlukan penanganan praproses yang berbeda. Koreksi pencahayaan pada kondisi
minim cahaya yang tidak seragam telah diteliti oleh \textcite{fu2026twolevel}.

\section{Pengenalan Karakter Optis untuk Deret Angka}

Pengenalan deret angka tanpa segmentasi karakter umumnya memakai arsitektur
CRNN yang diperkenalkan \textcite{shi2017crnn}. Arsitektur tersebut menggabungkan
lapisan konvolusional untuk mengekstraksi ciri, lapisan \textit{bidirectional
long short-term memory} untuk memodelkan urutan, dan fungsi kerugian CTC untuk
menyelaraskan urutan keluaran dengan label tanpa segmentasi eksplisit
\parencite{graves2006ctc}.

CTC memiliki keterbatasan struktural yang relevan bagi penelitian ini. Karena
segmentasi karakter dianggap tidak diketahui, dua karakter identik yang
berurutan hanya dapat dihasilkan apabila terdapat simbol kosong di antaranya.
Akibatnya deret seperti ``1006.5'' berpotensi terbaca ``106.5''. Pola kegagalan
tersebut terdokumentasi pada implementasi CRNN untuk tampilan tujuh segmen
\parencite{oicws2026lcd}. Alternatif dekoder tanpa rekurensi dikembangkan
melalui NRTR \parencite{sheng2019nrtr} dan PARSeq
\parencite{bautista2022parseq}.

\section{PP-OCRv6}

PP-OCRv6 adalah keluarga model pengenalan karakter yang dirilis pada Juni 2026
\parencite{zhang2026ppocrv6}. Model ini tersedia dalam tiga tingkat, yaitu
\textit{tiny} berukuran 1,5 juta parameter, \textit{small} berukuran 7,7 juta
parameter, dan \textit{medium} berukuran 34,5 juta parameter. Pada tingkat
\textit{small}, kepala pengenalannya berukuran sekitar 5,2 juta parameter,
sehingga angka tersebut tidak sama dengan ukuran keseluruhan model.

Secara arsitektural, PP-OCRv6 memakai tulang punggung PPLCNetV4 yang bergaya
\textit{MetaFormer} dengan reparameterisasi struktural, leher deteksi RepLKFPN,
serta leher pengenalan \textit{EncoderWithLightSVTR} yang memakai atensi
lokal-global. Pengenalan dilakukan oleh dekoder gabungan CTC dan NRTR, dengan
hanya cabang CTC yang dievaluasi saat inferensi; cabang NRTR berperan pada saat
pelatihan sebagai pengatur tambahan bagi representasi enkoder
\parencite{zhang2026ppocrv6}. Model ini tersedia melalui Paddle, Transformers,
dan ONNX Runtime, sehingga dapat dijalankan tanpa koneksi jaringan apabila
bobotnya dibundel bersama aplikasi. Model generasi sebelumnya, PP-OCRv5,
menunjukkan bahwa model berukuran 5 juta parameter dapat menyaingi model
bahasa-visi berukuran miliaran parameter pada tugas pengenalan karakter
\parencite{zhang2026ppocrv5}.

\section{Model Bahasa-Visi}

Model bahasa-visi menerima citra dan pertanyaan dalam bentuk teks, lalu
menghasilkan jawaban. Pada domain alat ukur, kemampuan model ini terbelah.
\textcite{ega2026leveraging} melaporkan \textit{full-string recognition} sebesar
93,85 persen dan akurasi per digit 98,10 persen untuk Gemini pada perangkat ukur
digital. Sebaliknya, \textcite{lin2026measurebench} melaporkan bahwa model
terbaik hanya mencapai 30,2 persen pada data nyata dari 26 jenis alat ukur.
Perbedaan tersebut dijelaskan oleh jenis tugasnya: transkripsi digit yang telah
tertulis versus inferensi nilai dari posisi jarum dan skala. Pada domain
penelitian ini, yaitu transkripsi digit, model bahasa-visi bukan pesaing
akurasi, melainkan pesaing dari sisi biaya komputasi dan kebutuhan jaringan.

\section{Keandalan dan \textit{Selective Prediction}}

Akurasi dan keandalan merupakan dua besaran yang berbeda. Akurasi mengukur
seberapa sering model benar, sedangkan keandalan menyangkut kemampuan sistem
menyatakan bahwa dirinya tidak layak dipercaya pada suatu kasus.
\textit{Selective prediction} memberi sistem kewenangan menolak memprediksi
apabila keyakinannya rendah \parencite{lin2026measurebench}. Pada kerangka ini,
kinerja tidak dinyatakan sebagai satu angka, melainkan sebagai hubungan antara
\textit{coverage}, yaitu proporsi kasus yang diterima otomatis, dan akurasi pada
kasus yang diterima. Hubungan tersebut digambarkan sebagai kurva
\textit{risk--coverage} dan diringkas sebagai luas daerah di bawah kurva atau
\textit{area under the risk--coverage curve} (AURC).

Nilai keyakinan mentah dari model tidak dapat langsung dipakai sebagai peluang
kebenaran. Diperlukan kalibrasi pada data yang terpisah dari data penilaian,
agar ambang keputusan tidak ditentukan dari data yang sama dengan data uji.

\section{Penelitian Terdahulu}

Ringkasan penelitian terdahulu yang menjadi acuan disajikan pada
Tabel~\ref{tab:terdahulu}.

\begin{table}[H]
\caption{Ringkasan penelitian terdahulu}
\label{tab:terdahulu}
\begin{tabularx}{\textwidth}{@{}p{0.6cm} p{3.3cm} p{4.6cm} X@{}}
\toprule
\textbf{No} & \textbf{Sumber} & \textbf{Hasil yang dilaporkan} & \textbf{Keterbatasan untuk penelitian ini} \\
\midrule
1 & \textcite{jiang2026method} & 99,7 persen pada 1.146 karakter dalam 20 milidetik, diuji pada timbangan elektronik & Bergantung pada kondisi citra; perilaku saat gagal tidak dilaporkan \\
2 & \textcite{effendi2025unjuk} & \textit{Exact match} 90,37 persen dengan CER 3,6 persen pada data logger waktu nyata & Berhenti pada akurasi; perilaku saat bacaan salah tidak dilaporkan \\
3 & \textcite{gushima2023automatic} & Generator citra sintetis dan model desimal terpisah, 96,8 persen & Terbatas pada perangkat berspesifikasi rendah; keandalan tidak diukur \\
4 & \textcite{yenneti2023lightweight} & 97,8 persen per bingkai pada 30 bingkai per detik, dengan penyaringan antarbingkai & Penyaringan dipakai untuk stabilisasi keluaran, bukan keputusan terima atau tolak \\
5 & \textcite{oicws2026lcd} & 98 persen pada data nyata, bobot terbuka & Angka terikat pada kumpulan data sendiri \\
6 & \textcite{low2024detecting} & DBNet terbaik untuk deteksi, PARSeq terbaik untuk pengenalan & Gagal pada kondisi pantulan cahaya \\
7 & \textcite{ega2026leveraging} & Gemini dengan \textit{full-string recognition} 93,85 persen & Membutuhkan jaringan dan biaya; keandalan tidak diukur \\
8 & \textcite{lin2026measurebench} & Model bahasa-visi terbaik hanya 30,2 persen pada data nyata & Menguji inferensi skala, bukan transkripsi digit \\
9 & \textcite{zhang2026ppocrv6} & PP-OCRv6, 1,5--34,5 juta parameter, 83,2 persen rata-rata 15 kategori & Belum diuji pada tampilan tujuh segmen \\
\bottomrule
\end{tabularx}
\end{table}

\section{Posisi Penelitian}

Berdasarkan kajian pada bagian sebelumnya, terdapat tiga celah yang belum
terjawab. Pertama, keandalan pembacaan tampilan tujuh segmen belum pernah
dilaporkan dalam bentuk laju kesalahan senyap, kurva \textit{risk--coverage},
maupun \textit{selective accuracy}. Kedua, belum ada penelitian yang memakai
konteks di luar citra, seperti jumlah slot digit, rentang berat yang sah untuk
suatu produk, toleransi penimbangan, dan konsistensi antarbingkai, sebagai dasar
keputusan kelayakan bacaan. Ketiga, PP-OCRv6 yang secara eksplisit menyebut
tampilan digital sebagai sasaran khusus belum diuji pada tampilan tujuh segmen
\parencite{zhang2026ppocrv6}.

Penelitian ini menempati celah tersebut dengan menggeser pertanyaan dari sekadar
``apakah bacaan dapat diperoleh'' menjadi ``kapan bacaan tidak layak dipercaya
dan bagaimana sistem mengetahuinya''. Pergeseran tersebut tidak bergantung pada
penggantian model, sehingga kontribusinya tidak dapat direduksi menjadi
pemakaian model yang lebih baru.
